\section{Background}

The emergence of Industry 4.0 has catalysed a paradigm shift in manufacturing operations, transforming conventional factories into cyber-physical production systems characterised by pervasive sensor instrumentation, real-time data acquisition, and autonomous decision-making. Smart manufacturing integrates industrial Internet of Things devices, cloud computing infrastructure, and advanced analytics to create self-optimising environments that respond dynamically to operational conditions. Within this landscape, maintenance strategies have evolved through four distinct paradigms: reactive maintenance, which addresses equipment failures after they occur, incurring high emergency costs and production losses; preventive maintenance, which applies periodic servicing according to fixed schedules regardless of actual machine condition, often resulting in unnecessary interventions and residual failure risk between service intervals; condition-based maintenance, which triggers actions upon detecting predefined degradation thresholds through continuous sensor monitoring; and predictive maintenance, which forecasts future failures from multivariate sensor trajectories to enable just-in-time intervention before functional breakdown. The economic imperative for this evolution is substantial, as unplanned downtime in high-throughput industrial settings carries severe financial consequences \cite{mobley_introduction_2002}. Jardine et al. \cite{jardine_review_2006} demonstrate that predictive maintenance programmes can reduce maintenance costs by 25--30\% and eliminate unplanned breakdowns by up to 75\% when properly implemented, establishing the compelling operational rationale for data-driven maintenance intelligence in modern smart factories.

In response to these challenges, machine learning and deep learning have emerged as the dominant methodological paradigms for anomaly detection and predictive maintenance in industrial time-series data. Chalapathy and Chawla \cite{Chalapathy2019} provide a comprehensive taxonomy of deep anomaly detection architectures, categorising autoencoders, recurrent neural networks, and convolutional networks across supervised, semi-supervised, and unsupervised learning settings. Pang et al. \cite{pang_deep_2021} extend this foundation with a systematic review demonstrating that reconstruction-based methods trained exclusively on normal data achieve superior performance in industrial scenarios with rare anomalous events. For multi-sensor environments, Gao et al. \cite{Gao2022} introduced TSMAE, a memory-augmented autoencoder that prevents anomalous inputs from being reconstructed with low error by retrieving representations from stored normal prototypes. Wen and Keyes \cite{wen_time_2019} proposed CNN-based segmentation approaches that outperform recurrent architectures on several benchmarks while achieving substantially lower computational cost. Han et al. \cite{Han2022} demonstrated that anomalous states can be detected by monitoring deviations from learned cross-sensor correlation structures, a finding directly applicable to the five-modality environment studied herein. These advances collectively establish that deep learning offers powerful tools for modelling complex temporal patterns in industrial sensor streams, yet the literature predominantly treats anomaly detection and remaining useful life estimation as isolated tasks rather than complementary components of a unified predictive maintenance workflow.

\section{Motivation}

Despite the theoretical advantages of predictive maintenance, its practical deployment in large-scale smart manufacturing facilities remains constrained by three systematic gaps in the current literature.

First, anomaly detection and remaining useful life estimation are predominantly treated as isolated research problems. Anomaly detection focuses on identifying outlier patterns in sensor streams \cite{chen_unsupervised_2020}, while prognostics research concentrates on forecasting degradation trajectories \cite{nizam_real-time_2022}, with limited integration into a unified decision workflow. Kendall et al. \cite{kendall_multi-task_2017} demonstrated that jointly learned representations improve individual task performance through shared encoding, yet this principle has not been systematically exploited in predictive maintenance systems.

Second, maintenance scheduling across large machine pools is a combinatorially hard problem that classical greedy heuristics only partially approximate \cite{dalyac_qualifying_2021}. As the number of machines grows, the solution space expands exponentially, and simple priority rules fail to capture pairwise conflict constraints and resource limitations inherent to real factory floors.

Third, static machine learning models degrade in accuracy as industrial environments evolve through sensor drift and shifting failure mode distributions \cite{lee_autoencoder-based_2024}. Without autonomous adaptation mechanisms, deployed systems require expensive manual retraining cycles, undermining the economic viability of predictive maintenance at scale.

These three gaps motivate the central research objective of this thesis: to design an integrated framework that unifies spatiotemporal sensor forecasting, multi-task supervised learning, autonomous multi-agent decision-making, and quantum-inspired combinatorial optimization within a single closed-loop architecture for smart manufacturing predictive maintenance.

\section{Contribution}

This thesis proposes the Quantum-Agentic Smart Manufacturing for Anomaly Detection and Predictive Maintenance (QASAMAP) framework as an integrated, real-time solution to the three limitations of Section~1.2. QASAMAP closes the loop from raw Kafka sensor streams to constraint-satisfying maintenance schedules within a single $60$-second cycle and is organised into three functional layers. \textbf{Layer~1} implements a Temporal Graph Convolutional Network (T-GCN) \cite{Zhao2020} for spatio-temporal sensor-stream forecasting, with a cross-correlation lag-graph over the 50 monitored machines to model machine-wide degradation propagation, plus a six-target supervised benchmark across XGBoost, LightGBM, CatBoost, HistGradientBoosting, and TabNet for anomaly, failure-type, downtime-risk, maintenance-required, machine-status, and remaining-useful-life prediction. \textbf{Layer~2} deploys a four-tier real-time multi-agent orchestrator: a Kafka consumer with a $W=10$ per-machine sliding window dispatches windows to nine specialised LLM-powered agents (Monitoring, Correlation, Pattern Learning, Diagnosis, Planning, Action, Self-Reflection, Explainability, Reporting) coordinated by a phased orchestrator with intra-phase parallelism, and emits a maintenance decision via the scheduling layer. Heterogeneous LLM routing (GLM-5.1, Kimi-K2.6, DeepSeek-v4-Pro) matches each agent to a task-appropriate model. \textbf{Layer~3} formulates maintenance scheduling as both Mixed-Integer Linear Programming (MILP) and Quadratic Unconstrained Binary Optimisation (QUBO), and benchmarks six solver paradigms on identical streaming-derived instances: classical exact MILP/CBC, the Quantum-Inspired Evolutionary Algorithm \cite{HanKim2002}, Simulated Quantum Annealing \cite{KadowakiNishimori1998}, feasibility-aware Simulated Annealing/Genetic Algorithm/Tabu Search, and the Quantum Approximate Optimisation Algorithm \cite{Farhi2014} with the XY ring mixer of Hadfield et al.\ \cite{Hadfield2019}, executed both in classical state-vector simulation and on IBM Heron~r2 hardware (\texttt{ibm\_fez}, 156 qubits) via the parameter-transfer strategy of Brandao et al.\ \cite{Brandao2018} and the iterative warm-start framework of Niroula et al.\ \cite{Niroula2026}. The XY-ring mixer preserves the one-hot machine-assignment constraint by construction, providing the constraint-native quantum path that the empirical results below validate. The overall three-layer architecture is illustrated in Figure~\ref{fig:research_methodology} and detailed in Chapter~3.

The primary contributions of this thesis are as follows.
\begin{itemize}
    \item \textbf{Spatio-temporal forecasting.} The T-GCN with cross-correlation lag-graph (50 nodes, 538 directed edges) achieves average $R^2 = 0.7702$ across five sensor channels on a $100{,}000$-reading industrial dataset, outperforming a tuned standalone LSTM ($R^2 = 0.6853$, $+12.4\%$) and a hybrid LSTM--GCN baseline ($R^2 = 0.7134$, $+8.0\%$).
    \item \textbf{Supervised health-prediction benchmark.} Five tree-based ensembles plus TabNet are evaluated across six machine-health targets, establishing reference performance for anomaly detection (best F1 $= 0.999$), failure-type classification (best F1 $= 0.880$), maintenance-required prediction (best F1 $= 0.872$), and remaining-useful-life regression (best $R^2 = 0.189$).
    \item \textbf{Real-time multi-agent orchestration.} The nine-agent phased pipeline closes the end-to-end loop from raw Kafka events to feasible maintenance schedule within the 60-s cycle. Intra-phase parallelism reduces worst-case latency from $9 \cdot T_{\text{LLM}}$ to $\sim 4 \cdot T_{\text{LLM}}^{\max}$. End-to-end validation on $50$ streamed machines forwards $45/50$ to the scheduler at CRITICAL priority and returns a one-hot- and capacity-feasible MILP schedule for all $45$ in $36.9$~ms.
    \item \textbf{Fair multi-paradigm scheduling benchmark.} A multi-seed scaling study on the basic constraint set ($N \in \{3,\ldots,30\}$, $5$ random seeds per $N$, identical $T$ and pure-cost metric) shows MILP/CBC dominates QIEA, SQA, and the QAOA simulator at every tested $N$, with sub-second solve times to $N=150$ — comfortably within the Kafka cycle.
    \item \textbf{Constraint-native quantum advantage on streaming-derived data.} Under the full complex-constraint set (one-hot, slot-window capacity, $\rho$-slot risk horizon, must-schedule), three measured regimes are identified on \texttt{ibm\_fez}: $N=3$ (depth $801$, $644$~CZ) where QAOA-XY post-selects the MILP optimum exactly at cost $300.00$ ($10.96\%$ feasibility); $N=5$ (depth $1{,}226$, $1{,}935$~CZ) where the problem is structurally infeasible for MILP and every classical heuristic, and QAOA-XY uniquely returns a feasible schedule at cost $1{,}186.02$ (only $1.3\%$ above MILP best-effort, $0.24\%$ feasibility); $N=7$ (depth $3{,}120$, $4{,}063$~CZ) where QAOA-XY beats every classical heuristic at cost $3{,}069.07$ with $V=2$ violations vs $V=8$ for the best classical at cost $4{,}295.42$.
    \item \textbf{Empirical NISQ noise frontier on Heron~r2.} The frontier sits at $N \approx 7$--$8$: $N=9$ (depth $2{,}958$, $5{,}390$~CZ) yields $0\%$ feasibility on \texttt{ibm\_fez}, with multibackend confirmation on \texttt{ibm\_kingston} and a parallel \texttt{ibm\_marrakesh} attempt. The limit on Heron~r2 at $p=2$ is set by gate fidelity and circuit depth, not qubit count \cite{Blekos2024,Preskill2018}.
\end{itemize}

By unifying spatio-temporal forecasting, agentic cognitive reasoning, and constraint-aware scheduling on a fair shared benchmark and a single 60-s real-time budget, QASAMAP provides both an honest reproducible benchmark and the empirical identification of the constraint-tight streaming-derived regime in which constraint-native quantum optimisation provides operational value over the best classical alternatives.

\section{Structure of Thesis}

The remainder of this thesis is organized as follows. Chapter 2 provides a review of the relevant literature across sensor-based anomaly detection, deep learning for predictive maintenance, quantum machine learning, and agentic AI systems. Chapter 3 presents the proposed QASAMAP methodology, including the T-GCN architecture, quantum computing layer design, and multi-agent orchestration system. Chapter 4 presents the data ingestion infrastructure and the spatiotemporal forecasting experiments using the Temporal Graph Convolutional Network, alongside the supervised classification and regression benchmarks. Chapter 5 details the multi-agent system for monitoring, diagnosis, planning, and action, including the LLM benchmark and post-audit comparative evaluation. Chapter 6 presents the quantum-inspired optimization formulations for maintenance scheduling, including Pauli-Z encoding, QUBO priority ranking, and QAOA graph scheduling. Chapter 7 concludes with a synthesis of contributions, limitations, and directions for future research.
